% Copyright (c) 2026 Geovana Neves. Licensed under CC BY 4.0 (https://creativecommons.org/licenses/by/4.0/). See LICENSE-AND-AUTHORSHIP.md.
%
% solvers/text/03-xfoil-qprop-choosing.tex
%
% Parts 5 to 7 of guide 3: XFOIL, QPROP and how to choose.

\section{XFOIL}

\begin{frame}[fragile]{XFOIL: an airfoil polar}
\begin{columns}[T]
\begin{column}{0.5\textwidth}
\begin{lstlisting}[style=cfdbtoml]
# inputs/setups/s001.toml
airfoil = "NACA 0012"
reynolds = 1000000
[solver]
launcher = ["wsl", "-d", "Ubuntu", "--"]
executable = "/home/you/bin/xfoil"
\end{lstlisting}
\begin{lstlisting}
# inputs/matrices/study.fs
POL  | RUN | SET | REF | PPROC | FLIGHT_CONDITION | SWEEP_VALUES
1001 | 1   | 001 | 001 | 001   | ALPHA:sweep      | 0,2,4
\end{lstlisting}
\end{column}
\begin{column}{0.46\textwidth}
\begin{itemize}\small
  \item \code{airfoil} is a NACA four or five digit name;
        \code{airfoil_file} is a coordinate file in the workspace instead.
  \item \code{reynolds} 0, the default, is inviscid. \code{mach} defaults
        to 0, \code{ncrit} to 9, \code{iterations} to 200.
  \item The r and p files may be empty.
\end{itemize}
\end{column}
\end{columns}
\begin{cfdbevidence}[Reading the result]
\code{CL}, \code{CD}, \code{CDp} and \code{CM} are kept as XFOIL prints them:
unit chord, moment about the quarter chord, nose up positive. A point is
\code{converged} with exactly one polar row at the requested angle, and
\code{failed} otherwise.
\end{cfdbevidence}
\end{frame}

\section{QPROP}

\begin{frame}[fragile]{QPROP: a propeller from its blade}
\begin{columns}[T]
\begin{column}{0.5\textwidth}
\begin{lstlisting}[style=cfdbtoml]
# inputs/setups/s001.toml
blade = "inputs/mesh/blade.obj"
n_blades = 3
n_sections = 12
spacing = "cosine"
cut = "cylinder"
solver_2d = "xfoil"

[solver]
launcher = ["wsl", "-d", "Ubuntu", "--"]
executable = "/home/you/bin/qprop"
xfoil_executable = "/home/you/bin/xfoil"
\end{lstlisting}
\end{column}
\begin{column}{0.46\textwidth}
\begin{itemize}\small
  \item \code{blade} is a 3D blade as an OBJ, in metres, or a set under
        \code{inputs/polars-2D/} with the stations and their polars. A blade
        named only in the \code{GEOMETRY} cell runs too.
  \item An OBJ is sliced at \code{n_sections} radial stations: radius,
        chord, pitch and section shape are measured.
  \item Each section gets its polar from XFOIL, at its own Reynolds and
        Mach number.
  \item \code{axis}, \code{hub_center} and \code{rotation} say how the
        blade sits in the file; the defaults are the x axis, the origin and
        counter-clockwise seen from the front.
\end{itemize}
\end{column}
\end{columns}
\end{frame}

\begin{frame}[fragile]{QPROP: the operating point and the result}
\begin{lstlisting}
POL  | RUN | SET | REF | PPROC | FLIGHT_CONDITION                              | SWEEP_VALUES
1001 | 1   | 001 | 001 | 001   | V:20, J:sweep, BETA07:30, RHO:1.2, MU:1.8e-5  | 0.6,0.8,1.0
\end{lstlisting}
\begin{columns}[T]
\begin{column}{0.48\textwidth}
\relax{\cfdbtablesize
\begin{tabularx}{\linewidth}{@{}P{0.36\linewidth}R@{}}
\toprule
\textbf{Key} & \textbf{Meaning} \\
\midrule
\code{V} & flight speed, m/s (required) \\
\code{RPM} or \code{J} & rotation rate; $J = V/(nD)$ \\
\code{BETA07} or \code{DBETA} & pitch at 0.7 R, or a pitch offset, degrees \\
\code{RHO}, \code{MU}, \code{A} & the air; sea level if left out \\
\bottomrule
\end{tabularx}}
\end{column}
\begin{column}{0.48\textwidth}
\relax{\cfdbtablesize
\begin{tabularx}{\linewidth}{@{}P{0.42\linewidth}R@{}}
\toprule
\textbf{Column} & \textbf{Meaning} \\
\midrule
\code{T}, \code{Q}, \code{P} & thrust (N), torque (N m), power (W) \\
\code{CT}, \code{CP} & $T/(\rho n^2 D^4)$, $P/(\rho n^3 D^5)$ \\
\code{CT_QPROP}, \code{CP_QPROP} & as QPROP prints them \\
\code{EFFICIENCY} & $TV/P$ \\
\bottomrule
\end{tabularx}}
\end{column}
\end{columns}
\vspace{1mm}
{\small\code{sections.json} in the point folder records each station and
where its polar came from. Read it when a result looks odd: a station whose
polar failed borrows its neighbour's, and the file says so.\par}
\end{frame}

\section{Choosing}

\begin{frame}{Which solver for which question}
\relax{\cfdbbodysize
\begin{tabularx}{\linewidth}{@{}P{0.52\linewidth}R@{}}
\toprule
\textbf{The question} & \textbf{Start with} \\
\midrule
The polar of an airfoil section, many angles, now & \code{xfoil} \\
Thrust and power of an isolated propeller over an advance-ratio sweep &
  \code{qprop} \\
A 2D section where separation or the turbulence model matters &
  \code{openfoam} or \code{su2} \\
A 3D configuration, compressible, with a volume mesh & \code{su2} \\
A 2D section whose unsteady flow must be resolved in time & \code{pyfr} \\
A 3D configuration without a volume mesh, or a rotating propeller &
  \code{flightstream}, through pyflightstream [8] \\
\bottomrule
\end{tabularx}}
\begin{takeaway}Start with the cheapest solver that resolves the effect you
are asking about.\end{takeaway}
\end{frame}

\begin{frame}{What no solver here does}
\begin{itemize}
  \item \textbf{Unsteady workflows.} Every library solver runs steady
        points. PyFR marches in time, and what it returns is a mean.
  \item \textbf{Mesh generation.} A mesh is an input. The package makes
        none, except the slicing of a propeller blade into sections.
  \item \textbf{A verdict on accuracy.} \code{converged} is about the run,
        not about the mesh or the model.
  \item \textbf{Mixing solvers in one workspace.} One workspace, one solver;
        compare two by giving both the same matrix.
\end{itemize}
\begin{cfdbevidence}[Comparing two solvers]
The matrix columns and the air resolved from a row are the same in every
workspace, and a point has the same folder name in each. The coefficient
names and conventions are not the same: Part 1 has the table.
\end{cfdbevidence}
\end{frame}
